% GENERATED FILE. Edit canonical lessons, then run npm run generate:books.
% canonical-source-hash: fnv1a64:8b7bcb355e488f86
% canonical-lessons: MR-C218-koli, MR-C218-mhais, MR-C218-bakri, MR-C218-dhuke, MR-C218-dav

\chapter{Spider, Buffalo, She-goat, Fog, Dew}
\label{ch:mr-spider-buffalo-she-goat-fog-dew}
\hlchaptermodality{218}

\begin{chapteropening}
\textbf{By the end of this chapter:} \emph{I can say a spider, a buffalo, a she-goat, fog and dew.}

Complete the last lesson of chapter 218: I can say a spider, a buffalo, a she-goat, fog and dew.
\end{chapteropening}

\section[\texorpdfstring{\mr{कोळी} (koḷī)}{कोळी (koḷī)}]{\mr{कोळी} (koḷī) — a spider}

\label{lesson:MR-C218-koli}



\begin{quote}
Before the new one: say the Marathi for a harbour, then the Marathi for a well.
\end{quote}

\subsection*{You'll want to know: \mr{कोळी}}
\textbf{\mr{कोळी}} — \emph{koḷī} — \textquotedblleft{}a spider\textquotedblright{}.

\begin{cousinweb}
\textbf{\mr{कोळी}} (\emph{koḷī}) is a tadbhava, a word inherited from Sanskrit and reshaped by sound change. It comes from \emph{kolika}, \textquotedblleft{}weaver\textquotedblright{} (CDIAL 3535). CDIAL glosses the Marathi word \textquotedblleft{}a caste of watercarriers, a sort of spider\textquotedblright{}. It cites a view linking the Sanskrit word with Tamil \emph{kōṭikar}, \textquotedblleft{}weaver\textquotedblright{}, and a Munda word for \textquotedblleft{}spider\textquotedblright{}.
\end{cousinweb}

\begin{grammarlens}[title={What you've built}]
One more word for this chapter.
\end{grammarlens}

\subsection*{Your turn}
\begin{itemize}
\raggedright
  \item \emph{Say it:} \emph{koḷī}
  \item \emph{Say it:} \emph{koḷī}, once more
  \item \emph{Say it:} say \emph{koḷī}
  \item \emph{Recall:} say the Marathi for a harbour, then the Marathi for a well, then say \emph{koḷī} again
\end{itemize}

\begin{tcolorbox}[breakable,colback=teal!4,colframe=teal!35!black,title={Before you move on}]
What does \mr{कोळी} mean? (\textquotedblleft{}A spider\textquotedblright{}.) Say it once more.
\end{tcolorbox}
\section[\texorpdfstring{\mr{म्हैस} (mhais)}{म्हैस (mhais)}]{\mr{म्हैस} (mhais) — a buffalo}

\label{lesson:MR-C218-mhais}



\begin{quote}
Before the new one: say the Marathi for a well, then the Marathi for a spider.
\end{quote}

\subsection*{You'll want to know: \mr{म्हैस}}
\textbf{\mr{म्हैस}} — \emph{mhais} — \textquotedblleft{}a buffalo\textquotedblright{}.

\begin{cousinweb}
\textbf{\mr{म्हैस}} (\emph{mhais}) is a tadbhava, a word inherited from Sanskrit and reshaped by sound change. It comes from \emph{mahiṣa}, \textquotedblleft{}great, powerful; buffalo\textquotedblright{} (CDIAL 9964). CDIAL cites H. W. Bailey: the buffalo was \textquotedblleft{}the great one of the beasts\textquotedblright{}.
\end{cousinweb}

\begin{grammarlens}[title={What you've built}]
One more word for this chapter.
\end{grammarlens}

\subsection*{Your turn}
\begin{itemize}
\raggedright
  \item \emph{Say it:} \emph{mhais}
  \item \emph{Say it:} \emph{mhais}, once more
  \item \emph{Say it:} say \emph{mhais}
  \item \emph{Recall:} say the Marathi for a well, then the Marathi for a spider, then say \emph{mhais} again
\end{itemize}

\begin{tcolorbox}[breakable,colback=teal!4,colframe=teal!35!black,title={Before you move on}]
What does \mr{म्हैस} mean? (\textquotedblleft{}A buffalo\textquotedblright{}.) Say it once more.
\end{tcolorbox}
\section[\texorpdfstring{\mr{बकरी} (bakrī)}{बकरी (bakrī)}]{\mr{बकरी} (bakrī) — a she-goat}

\label{lesson:MR-C218-bakri}



\begin{quote}
Before the new one: say the Marathi for a spider, then the Marathi for a buffalo.
\end{quote}

\subsection*{You'll want to know: \mr{बकरी}}
\textbf{\mr{बकरी}} — \emph{bakrī} — \textquotedblleft{}a she-goat\textquotedblright{}.

\begin{cousinweb}
\textbf{\mr{बकरी}} (\emph{bakrī}) is a tadbhava, a word inherited from Sanskrit and reshaped by sound change. It comes from \emph{barkara}, \textquotedblleft{}kid, lamb; goat\textquotedblright{} (CDIAL 9153).
\end{cousinweb}

\begin{grammarlens}[title={What you've built}]
One more word for this chapter.
\end{grammarlens}

\subsection*{Your turn}
\begin{itemize}
\raggedright
  \item \emph{Say it:} \emph{bakrī}
  \item \emph{Say it:} \emph{bakrī}, once more
  \item \emph{Say it:} say \emph{bakrī}
  \item \emph{Recall:} say the Marathi for a spider, then the Marathi for a buffalo, then say \emph{bakrī} again
\end{itemize}

\begin{tcolorbox}[breakable,colback=teal!4,colframe=teal!35!black,title={Before you move on}]
What does \mr{बकरी} mean? (\textquotedblleft{}A she-goat\textquotedblright{}.) Say it once more.
\end{tcolorbox}
\section[\texorpdfstring{\mr{धुके} (dhuke)}{धुके (dhuke)}]{\mr{धुके} (dhuke) — fog}

\label{lesson:MR-C218-dhuke}



\begin{quote}
Before the new one: say the Marathi for a buffalo, then the Marathi for a she-goat.
\end{quote}

\subsection*{You'll want to know: \mr{धुके}}
\textbf{\mr{धुके}} — \emph{dhuke} — \textquotedblleft{}fog\textquotedblright{}.

\begin{cousinweb}
The origin of \textbf{\mr{धुके}} (\emph{dhuke}) is not settled. CDIAL lists a Marathi \emph{dhumẽ}, \textquotedblleft{}fog\textquotedblright{}, under Sanskrit \emph{dhūmra}, \textquotedblleft{}smoke-coloured\textquotedblright{} (CDIAL 6861), but has no entry for \textbf{\mr{धुके}}.
\end{cousinweb}

\begin{grammarlens}[title={What you've built}]
One more word for this chapter.
\end{grammarlens}

\subsection*{Your turn}
\begin{itemize}
\raggedright
  \item \emph{Say it:} \emph{dhuke}
  \item \emph{Say it:} \emph{dhuke}, once more
  \item \emph{Say it:} say \emph{dhuke}
  \item \emph{Recall:} say the Marathi for a buffalo, then the Marathi for a she-goat, then say \emph{dhuke} again
\end{itemize}

\begin{tcolorbox}[breakable,colback=teal!4,colframe=teal!35!black,title={Before you move on}]
What does \mr{धुके} mean? (\textquotedblleft{}Fog\textquotedblright{}.) Say it once more.
\end{tcolorbox}
\section[\texorpdfstring{\mr{दव} (dav)}{दव (dav)}]{\mr{दव} (dav) — dew}

\label{lesson:MR-C218-dav}



\begin{quote}
Before the new one: say the Marathi for a she-goat, then the Marathi for fog.
\end{quote}

\subsection*{You'll want to know: \mr{दव}}
\textbf{\mr{दव}} — \emph{dav} — \textquotedblleft{}dew\textquotedblright{}.

\begin{cousinweb}
\textbf{\mr{दव}} (\emph{dav}) is a tadbhava, a word inherited from Sanskrit and reshaped by sound change. It comes from \emph{drava}, \textquotedblleft{}running, flowing; fluidity\textquotedblright{} (CDIAL 6623). CDIAL glosses the Marathi word \textquotedblleft{}dew, dampness\textquotedblright{}. Bloch derived it instead from \emph{udaka}, \textquotedblleft{}water\textquotedblright{}.
\end{cousinweb}

\begin{grammarlens}[title={What you've built}]
One more word for this chapter.
\end{grammarlens}

\subsection*{Your turn}
\begin{itemize}
\raggedright
  \item \emph{Say it:} \emph{dav}
  \item \emph{Say it:} \emph{dav}, once more
  \item \emph{Say it:} say \emph{dav}
  \item \emph{Recall:} say the Marathi for a she-goat, then the Marathi for fog, then say \emph{dav} again
\end{itemize}

\begin{tcolorbox}[breakable,colback=teal!4,colframe=teal!35!black,title={Before you move on}]
What does \mr{दव} mean? (\textquotedblleft{}Dew\textquotedblright{}.) Say it once more.
\end{tcolorbox}
